\subsection{贪心}
贪心每一步选择当前最合适的方案，并且不回退。要说明这种选择能保留某个最优解；局部看起来更好，并不自动保证全局最优。
\subsubsection{按结束时间选区间}
从若干半开区间$[l_i,r_i)$中选最多个互不重叠的区间，允许一个区间在另一个结束时开始。按结束时间递增，选择每个与已选区间不冲突的候选。
% snippet:greedy
\begin{lstlisting}
struct Interval { int l, r; };
bool by_end(const Interval &a, const Interval &b) {
    return a.r < b.r;
}
// a[0..n-1], each interval satisfies l < r.
std::sort(a, a+n, by_end);
int count = 0, last = 0;
bool selected = false;
for (int i = 0; i < n; i++) {
    if (!selected || a[i].l >= last) {
        count++; last = a[i].r; selected = true;
    }
}
\end{lstlisting}
把某个最优方案中的第一个区间换成结束最早的区间，不会占用更多后续时间，因而不减少可选数量。对后续区间重复这个交换理由，得到贪心正确性。排序$O(n\log n)$，扫描$O(n)$，存储$O(n)$。
\subsubsection{中位数与距离和}
给定n个数，选择x最小化$\sum_{i=1}^n|a_i-x|$。排序后将最小与最大、次小与次大配对；每对的距离和在两者之间恒定，超过这个范围就增大。所有配对的最佳范围交于中位数处，所以任选一个中位数即可。
% snippet:median
\begin{lstlisting}
std::sort(a, a+n);
long long x = a[n/2], answer = 0;
for (int i = 0; i < n; i++) {
    long long d = (long long)a[i] - x;
    answer += d >= 0 ? d : -d;
}
\end{lstlisting}
n为偶数时，两项中位数之间的任意x都最优。时间$O(n\log n)$，存储数组$O(n)$。

\subsection{基础数学}
\subsubsection{最大公约数}
$\gcd(a,b)=\gcd(b,a\bmod b)$：两组数的公因数相同。反复取余，余数变为0时，剩余非零数就是最大公约数。
% snippet:gcd
\begin{lstlisting}
long long gcd(long long a, long long b) { // a, b >= 0
    while (b != 0) {
        long long r = a % b;
        a = b; b = r;
    }
    return a;
}
\end{lstlisting}
至少一个数非零时，时间$O(\log\max(a,b))$，空间$O(1)$。正整数的最小公倍数为$a/\gcd(a,b)\times b$；先除再乘减小中间值，但最终结果仍须在long long范围内。
\subsubsection{质因数分解}
正整数$x>1$可以写为$\prod p_i^{e_i}$，其中$p_i$是不同质数，$e_i$是对应指数。试除到$\sqrt{x}$，找到因子后反复整除。
% snippet:factor
\begin{lstlisting}
for (long long p = 2; p <= x / p; p++) {
    if (x % p != 0) continue;
    int exponent = 0;
    while (x % p == 0) { x /= p; exponent++; }
    printf("%lld %d\n", p, exponent);
}
if (x > 1) printf("%lld 1\n", x);
\end{lstlisting}
每次到达p时，更小的质因子已全部除尽，所以能整除的p一定是质数。结束后若$x>1$，它只能是一个剩余质因子。用$p\le x/p$避免$p^2$溢出。最坏时间$O(\sqrt{x_{\mathrm{初始}}})$，额外空间$O(1)$。
\subsubsection{矩阵乘法与结合顺序}
A为$n\times k$，B为$k\times m$时，$C=AB$是$n\times m$矩阵，
\[
C_{i,j}=\sum_{t=0}^{k-1}A_{i,t}B_{t,j}.
\]
% snippet:matrix
\begin{lstlisting}
for (int i = 0; i < n; i++)
    for (int j = 0; j < m; j++) {
        c[i][j] = 0;
        for (int t = 0; t < k; t++)
            c[i][j] += 1LL * a[i][t] * b[t][j];
    }
\end{lstlisting}
c使用long long且需能容纳累加结果。时间$O(nkm)$，结果空间$O(nm)$。矩阵乘法满足结合律：链式乘法可以改变括号，先计算产生较小中间矩阵的一对，但不能交换矩阵顺序。
